% Cluster A: notation and classical background
% Body-only LaTeX; preamble macros, cite keys, and cross-cluster labels per CONVENTIONS.md.

\section{Notation and classical background}\label{sec:prelim}

Everything in this section is classical.  It is collected once, in the
normalized forms used by the rest of the paper, for three reasons: the
statements are invoked repeatedly and in exact form; two of them --- the
compacta form of Alexander duality and the tameness of plane arcs --- are
easy to quote in weaker versions that would not support the proofs; and the
source manuscript's terminology --- \emph{atom}, \emph{package},
\emph{carrier}, \emph{fullification}, \emph{protected access},
\emph{collar} --- is fixed here at first use.

\subsection{Plane, sphere, and general notation}\label{subsec:A-notation}

Throughout, $\C$ is the complex plane, $\Z$ the integers, and
$\Chat=\C\cup\{\infty\}$ the Riemann sphere, identified with the round
sphere $S^2$ by stereographic projection and carrying its standard
orientation.  Compact subsets of $\C$ are exactly the compact subsets of
$\Chat$ missing $\infty$, and we pass between the two ambient spaces
without comment.  The symbol $\dist$ denotes Euclidean distance in $\C$;
diameters of crosscuts (Section~\ref{subsec:A-primeends}) are measured in
the spherical metric, and all other metric quantities are Euclidean.  For
a set $X$, $\Int X$ is its interior, $\partial X$ its boundary, and
$\overline{X}$ its closure, taken in the ambient space clear from context;
the relative closure of $X$ in a subspace $Y$ is written
$\overline{X}^{\,Y}$, as in the attachment trace \eqref{eq:B-trace}.

For sets $A,B$ we write $A\Subset B$, and say $A$ is \emph{compactly
contained} in $B$, if $\overline{A}$ is compact and
$\overline{A}\subset\Int B$; this is the manuscript's relation~$\Subset$.
A \emph{domain} is a nonempty connected open subset of $\C$ or $\Chat$; a
\emph{plane domain} is a domain in $\C$.  Every plane domain is polygonally
connected: any two of its points are joined by a polygonal path inside it
\cite{New51}.  A \emph{Jordan curve} is a homeomorphic image of the circle,
a \emph{Jordan arc} a homeomorphic image of $[0,1]$, and an \emph{open
Jordan arc} a Jordan arc with its two endpoints removed.  A \emph{Jordan
disk} (synonymously, a \emph{Jordan domain}) is a bounded plane domain
whose boundary is a Jordan curve; by the Jordan and Schoenflies theorems
recalled in Proposition~\ref{prop:A-extension} below, a Jordan disk is
simply connected, and its closure --- a \emph{closed Jordan disk} --- is
homeomorphic to the closed unit disk.  A \emph{continuum} is a nonempty
compact connected subset of $\C$ or $\Chat$; it is \emph{nondegenerate} if
it contains more than one point.

\subsection{Full continua, filled hulls, and fullification}
\label{subsec:A-full}

A compact set $K\subset\C$ is \emph{full} if $\C\setminus K$ is connected.
Following the source manuscript (equation~(9.5) of \cite{MWman}), the
\emph{filled hull} of a compact set $K\subset\C$ is
\begin{equation}\label{eq:A-fill}
\Fill(K)\;=\;K\,\cup\,
\bigcup\bigl\{\,B : B \text{ a bounded component of }\C\setminus K\,\bigr\},
\end{equation}
that is, $K$ together with the bounded components of its complement.  The
passage from $K$ to $\Fill(K)$ is called \emph{fullification}; the term is
the manuscript's, and it is used as a noun for the operation and,
occasionally, for its result.  The following properties are used
constantly and are recorded with proof.

\begin{lemma}[Basic properties of filled hulls]\label{lem:A-full}
Let $K\subset\C$ be compact, and let $\Omega_\infty(K)$ denote the
unbounded component of $\C\setminus K$.
\begin{enumerate}
\item $\Fill(K)=\C\setminus\Omega_\infty(K)$; in particular $\Fill(K)$ is
compact, and it is contained in every closed disk containing $K$.
\item The following are equivalent: $K$ is full; $\Chat\setminus K$ is
connected; $\Fill(K)=K$.
\item If $K$ is a continuum, then $\Fill(K)$ is a full continuum, and
$\Chat\setminus\Fill(K)=\Omega_\infty(K)\cup\{\infty\}$ is a simply
connected domain containing $\infty$.
\item The operator $\Fill$ is extensive ($K\subset\Fill(K)$), monotone
($K\subset L$ implies $\Fill(K)\subset\Fill(L)$), and idempotent
($\Fill(\Fill(K))=\Fill(K)$).
\end{enumerate}
\end{lemma}

\begin{proof}
(1) The set $\C\setminus\Omega_\infty(K)$ consists of $K$ together with
all bounded components of $\C\setminus K$, which is \eqref{eq:A-fill}; it
is closed because $\Omega_\infty(K)$ is open.  If $E$ is a closed disk with
$K\subset E$, then $\C\setminus E$ is a connected unbounded subset of
$\C\setminus K$, hence contained in $\Omega_\infty(K)$; taking complements
gives $\Fill(K)\subset E$.  Thus $\Fill(K)$ is closed and
bounded.

(2) The components of $\Chat\setminus K$ are the bounded components of
$\C\setminus K$ together with $\Omega_\infty(K)\cup\{\infty\}$: the latter
set is open in $\Chat$ (a spherical neighborhood of $\infty$ misses the
compact set $K$, and its intersection with $\C$, being connected and
unbounded, lies in $\Omega_\infty(K)$), it is connected, and its
complement within $\Chat\setminus K$ is the union of the bounded
components, which is open.
Hence $\Chat\setminus K$ is connected exactly when there are no bounded
components, which is fullness of $K$, which by \eqref{eq:A-fill} is the
identity $\Fill(K)=K$.

(3) Let $K$ be a continuum and let $B$ be a bounded component of
$\C\setminus K$.  The boundary of $B$ is contained in $K$: if a boundary
point $x$ of $B$ lay in $\C\setminus K$, its component $B'$ in
$\C\setminus K$ would be an open set containing $x$, hence meeting $B$,
hence equal to $B$; but then $x\in B$, contradicting openness of $B$ at
its own boundary point.  So $\partial B\subset K$.  Thus $\overline{B}=B\cup\partial B$ is a
connected set meeting $K$, and
$\Fill(K)=K\cup\bigcup_B\overline{B}$ is a union of connected sets each
meeting the connected set $K$, hence connected.  With (1) this makes
$\Fill(K)$ a continuum, and its complement in $\Chat$ is the domain
$\Omega_\infty(K)\cup\{\infty\}$, so $\Fill(K)$ is full.  Finally, a
domain in $\Chat$ whose complement is connected is simply connected
\cite{New51,Why42}, and the complement of the domain
$\Omega_\infty(K)\cup\{\infty\}$ is the continuum $\Fill(K)$.

(4) Extensivity is \eqref{eq:A-fill}.  For monotonicity, let
$K\subset L$; then $\Omega_\infty(L)$ is a connected unbounded subset of
$\C\setminus K$, hence $\Omega_\infty(L)\subset\Omega_\infty(K)$, and
taking complements in (1) gives $\Fill(K)\subset\Fill(L)$.  For
idempotence, the complement of $\Fill(K)$ is the single unbounded
component $\Omega_\infty(K)$, so $\Fill(\Fill(K))=\C\setminus
\Omega_\infty(K)=\Fill(K)$.
\end{proof}

The deeper absorption law
$\Fill(\Fill(K)\cup L)=\Fill(K\cup L)$, which makes fullification of a
typed packet independent of processing order, is
Lemma~\ref{lem:fill-absorption} below (manuscript Lemma~10.12); it is a
formal consequence of the three properties in clause~(4).

\begin{definition}[Package, atoms, carrier, fullification]
\label{def:A-package}
A \emph{package} is a nonempty finite union
$S=A_1\cup\dots\cup A_s$, $s\ge1$, of pairwise disjoint full continua,
called its \emph{atoms}.  Atoms may be
degenerate: in Theorem~\ref{thm:no-least} the package consists of the two
points $\pm1$.  A \emph{carrier} of the package $S$ is a compact
connected full continuum $H\subset\C$ with $S\subset H$.  Both
\emph{atom} --- a required full continuum of the package --- and
\emph{fullification} --- the passage \eqref{eq:A-fill} --- are the source
manuscript's terms \cite{MWman} and are used in its sense throughout.
\end{definition}

\subsection{Janiszewski's separation theorem}\label{subsec:A-janiszewski}

A closed set $F\subset\Chat$ \emph{separates} points
$p,q\in\Chat\setminus F$ if $p$ and $q$ lie in different components of
$\Chat\setminus F$; the set $F$ is \emph{nonseparating} if
$\Chat\setminus F$ is connected.  By Lemma~\ref{lem:A-full}(2), a compact
subset of $\C$ is nonseparating in $\Chat$ exactly when it is full; the
two words are used interchangeably for such sets.

\begin{theorem}[Janiszewski's separation theorem;
{\cite{Jan15}; see also \cite{New51,Why42}}]
\label{thm:A-janiszewski}
Let $F_1,F_2\subset\Chat$ be closed sets whose intersection
$F_1\cap F_2$ is connected and nonempty, and let
$p,q\in\Chat\setminus(F_1\cup F_2)$.  If neither $F_1$ nor $F_2$
separates $p$ from $q$, then $F_1\cup F_2$ does not separate $p$ from
$q$.
\end{theorem}

\begin{corollary}\label{cor:A-janiszewski}
If $F_1,F_2\subset\Chat$ are closed and nonseparating and
$F_1\cap F_2$ is connected and nonempty, then $F_1\cup F_2$ is
nonseparating.  In particular, the union of two full compact subsets of
$\C$ whose intersection is connected and nonempty is full.
\end{corollary}

\begin{proof}
Let $p,q\in\Chat\setminus(F_1\cup F_2)$.  Since $\Chat\setminus F_1$ is
connected, $F_1$ does not separate $p$ from $q$, and likewise for $F_2$.
By Theorem~\ref{thm:A-janiszewski}, $F_1\cup F_2$ does not separate $p$
from $q$.  As the pair was arbitrary, $\Chat\setminus(F_1\cup F_2)$ is
connected.  The final sentence follows from Lemma~\ref{lem:A-full}(2).
\end{proof}

This is the form used in the Janiszewski induction of
Lemma~\ref{lem:tree-full}, where the intersections are singletons.  The
complementary situation of finitely many pairwise \emph{disjoint} full
compacta, where the stated form of the theorem does not apply, is handled
cohomologically, by the duality of Section~\ref{subsec:A-duality} and the
nerve argument of Lemma~\ref{lem:packets} (manuscript Lemma~10.10).

\subsection{Prime ends, crosscut chains, and protected accesses}
\label{subsec:A-primeends}

Informally, an \emph{access} to a compact set $K$ is a way of reaching
$\partial K$ from the complement.  It can be presented by a landing arc,
as the access $p_0^{+}(m)$ of Proposition~\ref{prop:access-collision} is
presented by the vertical segment $m+it$; but landing presentations
presuppose that something lands, and for the arbitrary continua handled in
this paper a landing presentation need not be available.  The presentation that requires
nothing is Carath\'eodory's theory of prime ends \cite{Pom92}, which we
now recall in the normalization used throughout.

Let $\Omega\subset\Chat$ be a simply connected domain whose complement
contains at least two points.  The standing example is
$\Omega=\Chat\setminus K$ for a nondegenerate full continuum $K$, in which
case $\Omega$ is a simply connected domain containing $\infty$ by
Lemma~\ref{lem:A-full}(3); Jordan disks are a second example.  A
\emph{crosscut} of $\Omega$ is an open Jordan arc $\gamma\subset\Omega$
whose closure $\overline{\gamma}$ is a Jordan arc with two distinct
endpoints on $\partial\Omega$.  A crosscut divides $\Omega$ into exactly
two components, each a simply connected domain having $\gamma$ on its
boundary \cite{Pom92}.

A \emph{crosscut chain} (or \emph{prime-end chain}; Pommerenke's null
chains \cite{Pom92}, in the following normalization) is a sequence
\[
\Gamma=(\gamma_n,D_n)_{n\ge1}
\]
in which the $\gamma_n$ are pairwise disjoint crosscuts of $\Omega$,
$\overline{\gamma_n}$ are pairwise disjoint,
$D_n$ is one of the two components of
$\Omega\setminus\gamma_n$, and
\begin{equation}\label{eq:A-null-chain}
 \overline{D_{n+1}}\cap\Omega\subset D_n,
 \qquad
 \operatorname{diam}_{\mathrm{sph}}(\gamma_n)\longrightarrow0 .
\end{equation}
The inclusion is the strict nesting condition, while diameter convergence
is the null condition.  The strict nesting implies both
$D_{n+1}\subset D_n$ and
$\gamma_{n+1}\subset\overline{D_{n+1}}\cap\Omega\subset D_n$; thus a
later crosscut lies strictly on the declared side of every earlier one.
This stronger formulation rules out the large-side pathologies that mere
nestedness of the open domains would permit.
The record of a chain thus contains strictly more than the crosscuts: it
contains, for every $n$, the declaration of the selected side, and it is
exactly this declaration that the certificates of
Section~\ref{sec:certificates} protect.  Two chains are
\emph{equivalent} if each eventually divides the other, that is, if every
selected domain of either chain contains all sufficiently deep selected
domains of the other.  An equivalence class is a \emph{prime end} of
$\Omega$.

Two classical facts about prime ends are used.  First, the
Carath\'eodory correspondence: an orientation-preserving Riemann map
$\psi\colon\{|z|<1\}\to\Omega$ induces a bijection between the points of
the unit circle and the prime ends of $\Omega$; radial approach converges
in the prime-end compactification, although a Euclidean radial limit need
not exist without additional boundary regularity; see
\cite[Chapter~2]{Pom92}.  Second, the
\emph{impression}
\[
I(\Gamma)\;=\;\bigcap_{n\ge1}\overline{D_n}
\]
(closures in $\Chat$) is a nonempty compact connected subset of
$\partial\Omega$ --- the intersection of a decreasing sequence of
nonempty compact connected sets --- and depends only on the prime end of
$\Gamma$, not
on the representing chain \cite{Pom92}.  The nested closures
$\overline{D_n}$ decrease \emph{to} the impression; this is the form in
which impressions enter later collar and exposure arguments.

The prime ends of $\Omega$ carry a natural \emph{cyclic order}: the one
transported from the unit circle by the Carath\'eodory correspondence.
It is independent of the choice of Riemann map, since two
orientation-preserving Riemann maps differ by a disk automorphism, which
preserves the cyclic orientation of the circle.  For a finite family of
prime ends represented by chains with pairwise disjoint selected sides,
the cyclic order can be computed inside the domain, with no reference to
a uniformization: choose one sufficiently deep oriented crosscut from
each chain; the cyclic order of the prime ends is the separation order of
these pairwise disjoint crosscuts in $\Omega$.  The computation is
carried out carefully, by means of embedded triods, in the proof of
Lemma~\ref{lem:crosscut-survival}, and that formulation is the one shown
to be invariant under carrier enlargement.

\begin{definition}[Protected access; collar]\label{def:A-protected}
Let $A$ be a nondegenerate atom.  A \emph{protected access} to $A$ is a
distinguished null prime-end crosscut chain
$\Gamma=(\gamma_n,D_n)_{n\ge1}$ of $\Chat\setminus A$, recorded together
with its selected sides.  The terminal selected domain $D_1$ is the
\emph{collar} of the access, and
\[
 \overline{D_1}\setminus A
\]
is its \emph{relative collar tail}.  When finitely many protected accesses
appear together inside a Jordan localization, their labels are required to
specify pairwise distinct prime ends at each fixed atom, and the chains are
\emph{truncated}---an initial segment of each chain is discarded and the
indices shifted---so that the relative terminal collar tails are pairwise
disjoint, lie in the localization, and miss the forbidden set and the other
atoms.  For chains based at one atom, disjoint sufficiently deep tails come
from Hausdorff separation in that atom's prime-end closed disk.  For chains
based at different atoms, and for avoidance of the forbidden compactum, the
same conclusion follows separately from compact separation together with
the fact that the nested compact sets
$\overline{D_{i,a,n}}$ decrease to the impression in $A_i$.
Their closures may meet inside a common atom: distinct prime ends can have
overlapping impressions, and no local connectivity assumption is being
smuggled into the normalization.  Section~\ref{sec:certificates} records
this convention as part of the certificate data.
\end{definition}

The adjective \emph{protected} names a discipline, not a topological
property: a protected access is one that every later construction is
forbidden to disturb, in the precise sense that no connector may enter
its relative collar tail (Definition~\ref{def:certificate}).  The reward for the
discipline is the survival theorem, Lemma~\ref{lem:crosscut-survival}:
the chain then persists verbatim, with the same crosscuts and the same
selected sides, as a prime-end chain of the enlarged complement.

\subsection{\v{C}ech cohomology and Alexander duality}
\label{subsec:A-duality}

For a compact set $K\subset S^2$, $\check H^*(K;\Z)$ denotes \v{C}ech
cohomology with integer coefficients.  Two of its properties are needed
\cite[Chapter~6]{Spa66}: continuity --- the \v{C}ech cohomology of $K$ is the direct
limit of the singular cohomology of the open neighborhoods of $K$, so
that every class and every relation is realized in some neighborhood ---
and the Mayer--Vietoris sequence for pairs of compact sets in $S^2$.
Both are used in the nerve computation of Lemma~\ref{lem:packets} and in
the second proof of Lemma~\ref{lem:tree-full}.  The duality statement is
the following.

\begin{theorem}[Alexander duality, compacta form;
{\cite[Chapter~6]{Spa66}}]\label{thm:A-duality}
For every nonempty proper compact subset $K\subset S^2$ and every integer
$q$, there is a natural isomorphism
\[
\widetilde{\check H}^{\,q}(K;\Z)
\;\cong\;
\widetilde H_{\,1-q}(S^2\setminus K;\Z).
\]
In particular,
\begin{equation}\label{eq:A-duality}
\widetilde H_0\bigl(S^2\setminus K;\Z\bigr)\;\cong\;\check H^1(K;\Z).
\end{equation}
For $q=2$, the same theorem gives $\check H^2(K;\Z)=0$, since
$S^2\setminus K$ is nonempty and its reduced homology in degree $-1$
vanishes.
\end{theorem}

The general statement above is the \v{C}ech-cohomology form for compacta
given by Spanier.  Hatcher's Corollary~3.45 gives the familiar locally
contractible special case \cite[pp.~255--257]{Hat02}.

This is equation~(9.6) of the manuscript \cite{MWman}.  Reading
\eqref{eq:A-duality} at the level of connectivity gives the dictionary
used throughout the paper.  The complement is open in a surface and hence
locally path connected, so its connected and path components agree and
singular $H_0$ detects connectedness.  Thus $S^2\setminus K$ is connected exactly when
$\check H^1(K;\Z)=0$, and hence, by Lemma~\ref{lem:A-full}(2), a compact
nonempty set $K\subset\C$ is full exactly when
$\check H^1(K;\Z)=0$.
``Nonseparating'' and ``vanishing first \v{C}ech cohomology'' are
synonyms, and the two proofs of Lemma~\ref{lem:tree-full} translate one
into the other.

\begin{warning}\label{warn:A-cech}
The \v{C}ech form of \eqref{eq:A-duality}, valid for an arbitrary
nonempty proper compact $K\subset S^2$, is required; the locally
contractible special case, in which \v{C}ech and singular cohomology
agree \cite{Hat02}, is not sufficient.  The atoms of this paper are
arbitrary full continua---in general not locally connected or locally
contractible, and not necessarily of the homotopy type of CW
complexes---and for such compacta duality genuinely fails with singular
cohomology in place of \v{C}ech.  Hatcher's standard quasi-circle (Warsaw
circle) $W\subset\C$, obtained by completing the closed topologist's sine
curve by an arc, is path connected and simply connected but has exactly two
complementary components \cite[Exercise~1.3.7 and pp.~255--256]{Hat02}.
Thus
$\widetilde H_0(S^2\setminus W;\Z)\cong\Z$ and compacta duality gives
$\check H^1(W;\Z)\cong\Z$, whereas $H_1(W;\Z)=0$ and the universal
coefficient theorem gives $H^1_{\mathrm{sing}}(W;\Z)=0$.  Every duality argument in this paper ---
including Lemma~\ref{lem:tree-full} (second proof),
Theorem~\ref{thm:separation-budget}, and Lemma~\ref{lem:packets} --- is therefore carried out in \v{C}ech
cohomology, and the nerve and continuity properties invoked there are the
\v{C}ech ones.
\end{warning}

\subsection{Jordan, Schoenflies, and Carath\'eodory extensions}
\label{subsec:A-extension}

The naturality statements of this paper
(Proposition~\ref{prop:naturality}, the naturality clauses of
Theorem~\ref{thm:audit}) rest on a package of classical extension
theorems, collected here with the exact attributes used later.

\begin{proposition}[Extension package]\label{prop:A-extension}
\leavevmode
\begin{enumerate}
\item \textup{(Jordan; Schoenflies)}  Every Jordan curve
$J\subset\Chat$ divides $\Chat$ into exactly two components, each a
simply connected domain with boundary $J$, and the closure of each is
homeomorphic to the closed unit disk.  Moreover, every homeomorphism
between two Jordan curves in $\Chat$ extends to a homeomorphism of
$\Chat$.
\item \textup{(Tameness of plane arcs)}  For every Jordan arc
$\alpha\subset\C$ there is a homeomorphism of $\C$ carrying $\alpha$ onto
the straight segment $[0,1]\times\{0\}$.  In particular every Jordan arc
extends to a Jordan curve, and no arc separates the plane.
\item \textup{(Carath\'eodory extension)}  Every conformal homeomorphism
$\phi\colon U\to V$ of Jordan disks extends to a homeomorphism
$\overline{U}\to\overline{V}$ of their closures.
\item \textup{(Glued sphere extension)}  Every orientation-preserving
conformal homeomorphism $\phi\colon U\to V$ of Jordan disks
$U,V\subset\C$ extends to an orientation-preserving homeomorphism
$\Phi\colon\Chat\to\Chat$, and $\Phi$ may be chosen with
$\Phi(\infty)=\infty$.
\item \textup{(Fill-equivariance)}  Every homeomorphism $\Phi$ of
$\Chat$ with $\Phi(\infty)=\infty$ satisfies
$\Phi(\Fill(K))=\Fill(\Phi(K))$ for every compact $K\subset\C$.
\end{enumerate}
\end{proposition}

\begin{proof}
Clauses (1) and (2) are classical plane topology; see
\cite{New51,Why42}.  Clause (3) is Carath\'eodory's theorem on boundary
extension of Riemann maps of Jordan domains; see \cite{Pom92}.

For (4), extend $\phi$ to a homeomorphism
$\overline{\phi}\colon\overline{U}\to\overline{V}$ by (3).  Its boundary
restriction is a homeomorphism $\partial U\to\partial V$ of Jordan
curves; the complementary sets $\Chat\setminus U$ and $\Chat\setminus V$
are closed Jordan disks by (1), and a homeomorphism between the boundary
circles of two closed disks extends over the disks (cone the
homeomorphism off an interior point, after identifying each closed disk
with the standard one by (1)).  Gluing the two extensions along the
boundary curves yields a homeomorphism $\Phi\colon\Chat\to\Chat$ with
$\Phi|_U=\phi$.  A sphere homeomorphism that preserves orientation on a
nonempty open set preserves it globally, so $\Phi$ is orientation
preserving.  Finally, $\infty$ lies in the interior of the complementary
disk $\Chat\setminus U$, and $\Phi(\infty)$ in the interior of
$\Chat\setminus V$; composing $\Phi$ with a homeomorphism of the closed
disk $\Chat\setminus V$ that fixes $\partial V$ pointwise and carries
$\Phi(\infty)$ to $\infty$ (such a homeomorphism exists, interiors of
disks being homogeneous relative to the boundary), extended by the
identity on $\overline{V}$, arranges $\Phi(\infty)=\infty$ without
disturbing $\Phi|_U=\phi$.

For (5), the restriction $\Phi|_{\C}$ is a homeomorphism of $\C$ onto
itself.  It maps $\C\setminus K$ homeomorphically onto
$\C\setminus\Phi(K)$ and permutes complementary components.  A
homeomorphism of $\C$ is proper, so the image of the unbounded component
of $\C\setminus K$ is unbounded and the images of bounded components are
bounded; hence $\Phi(\Omega_\infty(K))=\Omega_\infty(\Phi(K))$, and
taking complements as in Lemma~\ref{lem:A-full}(1) gives
$\Phi(\Fill(K))=\Fill(\Phi(K))$.
\end{proof}

Clause (4) is the mechanism behind every naturality statement in the
paper: a univalent change of Jordan localization extends to a sphere
homeomorphism, and a sphere homeomorphism preserves all the notions the
certificates are built from --- compactness, connectedness, fullness (by
clause (5) and Lemma~\ref{lem:A-full}), crosscuts, selected sides, null
chains, intersections and disjointness, and, when orientation is
preserved, the cyclic order of prime ends and the ribbon structure of an
embedded tree.  Clause (2) is used in the slit-straightening steps of the
repaired slit lemma (Lemma~\ref{lem:slits}).

This completes the classical inventory.  The dynamical vocabulary of the
source manuscript --- frozen inventory, free domain, ports, packets,
audits --- is introduced where it is first needed:
the frozen free domain $\Omfree$ in Section~\ref{sec:criterion}, packets
and their audit in Sections~\ref{sec:packets} and~\ref{sec:audit}.
